\chapter*{Preface}
\addcontentsline{toc}{chapter}{Preface}
\markboth{Preface}{}

\section*{Who this handbook is for}

This handbook is written for someone who has \emph{never} touched Kubernetes.
You do not need to know what a container is, what ``orchestration'' means, or
why everyone keeps talking about YAML. If you can open a terminal, run a
command, and read a little code, you have everything you need.

By the last page you will be able to:

\begin{itemize}
  \item explain, in plain words, what Kubernetes is and why it exists;
  \item describe every major component of a cluster and what it does;
  \item write and apply manifests for Pods, Deployments, Services, Gateways,
        ConfigMaps, Secrets, volumes, Jobs, StatefulSets and more;
  \item secure, scale, observe, and troubleshoot real workloads;
  \item package applications with Helm and Kustomize, extend Kubernetes with
        CRDs and Operators, and deliver changes with CI/CD and GitOps;
  \item take a real multi-service application from source code all the way to a
        monitored, autoscaled, GitOps-managed deployment.
\end{itemize}

\section*{How the book is organised}

The book climbs a staircase. Each part assumes only what came before it.

\begin{center}
\resizebox{\linewidth}{!}{%
\begin{tikzpicture}[
  stage/.style={regular polygon, regular polygon sides=7, fill=k8sblue,
               text=white, minimum size=2.1cm, inner sep=0pt, font=\bfseries\Large},
  lbl/.style={text width=2.6cm, align=center, font=\footnotesize\bfseries, text=k8snavy}]
  \foreach \rom/\lab/\col [count=\n] in {I/Foundations/k8sblue, II/Workloads/k8sblue,
       III/Networking/k8sdeep, IV/Config \& Storage/k8sdeep, V/Operations/k8snavy,
       VI/Advanced/k8snavy, VII/Capstone/k8sorange}{
    \node[stage, fill=\col] (s\n) at (\n*2.9, \n*0.45) {\rom};
    \node[lbl, below=0.15cm of s\n] {\lab};
  }
  \foreach \a/\b in {1/2,2/3,3/4,4/5,5/6,6/7}
    \draw[-{Stealth[length=3mm]}, line width=1.4pt, k8smid] (s\a) -- (s\b);
  \node[font=\footnotesize\itshape, text=k8sgray] at (2.9,-1.6) {beginner};
  \node[font=\footnotesize\itshape, text=k8sgray] at (20.3,1.6) {production-ready};
\end{tikzpicture}}
\end{center}

\begin{description}[leftmargin=0pt, itemsep=0.4em]
  \item[\textcolor{k8sblue}{Part I -- Foundations.}] Containers, the problem
        Kubernetes solves, cluster architecture, installing a local cluster, and
        the \kc{kubectl}/YAML workflow.
  \item[\textcolor{k8sblue}{Part II -- Running Workloads.}] Pods, labels and
        namespaces, Deployments and rollouts, StatefulSets, DaemonSets, Jobs and
        CronJobs.
  \item[\textcolor{k8sblue}{Part III -- Networking.}] Services, DNS, Ingress,
        the Gateway API, and NetworkPolicies.
  \item[\textcolor{k8sblue}{Part IV -- Configuration \& Storage.}] ConfigMaps,
        Secrets, Volumes, PersistentVolumes and StorageClasses.
  \item[\textcolor{k8sblue}{Part V -- Operating Kubernetes.}] Requests and
        limits, scheduling, autoscaling, security, logging, metrics and tracing.
  \item[\textcolor{k8sblue}{Part VI -- Advanced Kubernetes.}] Helm, Kustomize,
        CRDs, Operators, GitOps, CI/CD, and running clusters in production.
  \item[\textcolor{k8sorange}{Part VII -- Capstone.}] We build \textbf{TaskFlow},
        a small but realistic three-tier application, and do \emph{all} of the
        DevOps work on Kubernetes: images, manifests, Kustomize overlays, a Helm
        chart, a CI pipeline, GitOps with Argo CD, autoscaling, network
        policies, monitoring and alerts.
\end{description}

\section*{How to read the boxes}

Throughout the book, coloured boxes highlight different kinds of information:

\begin{keyidea}
A concept you must remember. If you only read the Key Idea boxes, you will
still understand the shape of Kubernetes.
\end{keyidea}

\begin{analogy}
An everyday comparison that makes an abstract idea concrete.
\end{analogy}

\begin{tryit}
A short hands-on exercise. Do them; Kubernetes is learned with your fingers.
\end{tryit}

\begin{warningbox}
A common mistake, a trap, or something that can hurt you in production.
\end{warningbox}

\begin{protip}
A habit or shortcut used by experienced engineers.
\end{protip}

\begin{deepdive}
Optional detail about how something works internally. Skip on first reading.
\end{deepdive}

\section*{Conventions}

Commands you type are shown in terminal boxes. Files are shown with their file
name in the title bar:

\begin{bashcode}{terminal}
# Lines starting with # are comments; you do not type them.
kubectl get pods
\end{bashcode}

\begin{yamlcode}{hello.yaml}
apiVersion: v1
kind: Pod
metadata:
  name: hello
\end{yamlcode}

Kubernetes moves quickly. Everything here targets Kubernetes \textbf{v1.30 and
newer}. Where a tool version appears (for example a Helm chart version), check
the project's site for the latest release before you copy it.

\section*{Trademarks}

Kubernetes and the Kubernetes logo are trademarks of The Linux Foundation. The
Kubernetes resource icons come from the official \kc{kubernetes/community}
repository. Docker, Helm, Prometheus, Argo, etcd, containerd, Istio and other
names and logos belong to their respective owners and are used here only to
identify those projects.
